\section{Discussion and Open Questions}
\label{sec:discussion}

\subsection{What this paper establishes}

Five results emerge from the analysis.

\paragraph{1.\ The phase boundary exists and is characterized.}
Compound feedback loops are not inherently catastrophic or inherently
self-correcting.  The outcome depends on the ratio of self-correction
rate to observation-loss rate, mediated by channel redundancy.
Theorem~\ref{thm:phase_boundary} gives the boundary; the linearized
sufficient condition (Equation~\ref{eq:self_correcting}) gives the
computable form.

\paragraph{2.\ All three feedback loops are the same dynamical system.}
The unified treatment (Section~\ref{sec:unified_loops}) shows that
structural avoidance, partial subsumption, and over-calibrated
restriction share a common formal structure.  Solutions for one loop
(controlled relaxation for the Wamura problem, structural discovery
incentives for leverage blindness) improve the phase boundary for all.

\paragraph{3.\ The Wamura problem is a degenerate phase boundary.}
Without a mechanism for evidence generation against unfalsifiable risk
claims, the entire state space is absorbing
(Proposition~\ref{prop:wamura}).  The Wamura problem was called ``the
most important open problem'' in the discussion of~\citet{lasser2026horizon}; the phase
boundary analysis explains why: it is the case where the
self-correction rate is zero by construction, making the feedback loop
absorbing regardless of all other parameters.

\paragraph{4.\ Cross-substrate redundancy is the load-bearing quantity.}
Same-substrate channels may share failure modes; cross-substrate
channels are structurally independent (under the channel
isolation property of~\citet{lasser2026exo} and, for additivity, the joint-factorization assumption).
The design criterion (Theorem~\ref{thm:redundancy}) is stated in terms
of cross-substrate redundancy, making it robust to correlated failures
within substrates.

\paragraph{5.\ In two-substrate systems, each substrate's channels are singularly load-bearing.}
This is not a normative claim about human importance.  It is a
structural prediction from the dynamics: in $m = 2$ systems, losing one
substrate's channel on a dimension reduces $\rhocross_k$ to at most
$1$, violating the redundancy criterion on any dimension where
$\Kmin(k, \Lyap_{\mathrm{avg}}) > 1$ and leaving zero margin where $\Kmin(k, \Lyap_{\mathrm{avg}}) \leq 1$
(Corollary~\ref{cor:two_substrate}).
The skeleton-substrate scenario (Section~\ref{sec:worked_example})
shows that progressive \emph{attenuation} of biological channels, not
outright removal, is the primary vulnerability.

\subsection{Limitations}

\paragraph{Linearized treatment.}
The phase boundary is stated in a linearized regime (constant $\rS$,
$\rW$, $\rsub$).  In the full dynamics these quantities are
state-dependent: a larger world-model error produces more aggressive
subsumption ($\rsub$ increases with $\Lyap$), and the correction rate
$\rS$ depends on $T_s$ which itself depends on the error history.  The
linearized treatment produces the qualitative result (self-correcting
and absorbing basins exist, separated by a characterized boundary) but
the quantitative location of the boundary shifts under the nonlinear
dynamics.  A full nonlinear analysis would tighten the boundary at the
cost of substantially more complex proofs.

\paragraph{Honest-channel assumption.}
The analysis assumes observation channels are honest: their signal is
noisy but not adversarially manipulated.  The cryptographic
commitment protocol of~\citet{lasser2026exo} provides the enforcement mechanism for this
assumption, and the two papers are designed to be read as complements:
the exogenous verification framework ensures channels are honest, this paper ensures enough honest
channels survive.  But the boundary between ``noisy'' and
``adversarial'' is itself a modeling choice, and in practice the line
is blurred by sophisticated manipulation strategies that operate within
the noise distribution.

\paragraph{Partial self-diagnosability.}
As noted in Section~\ref{sec:design_criterion}, the actor can observe
$\rhocross_k$ and $\rS$ but not $\rW$ or $\dmax$.  The phase boundary
is therefore a sufficient condition for self-correction that the actor
can partially monitor, not a decision procedure it can fully evaluate.
This is an honest limitation, not a fixable gap: the quantities the
actor cannot observe are precisely the ones that depend on the true
dynamics rather than the perceived dynamics.

\subsection{Connections to other proposals}

\paragraph{Controlled relaxation.}
Proposition~\ref{prop:wamura} shows the Wamura problem is a degenerate
phase boundary.  Controlled relaxation protocols (the subject of a
companion paper) provide the mechanism that makes $\rS > 0$ for
Loop~3.  The two analyses are complementary: this paper provides the
dynamical framework showing \emph{why} a positive self-correction rate
matters, while controlled relaxation provides the specific protocol
that \emph{achieves} it.

\paragraph{The B-to-C gap.}
The phase boundary analysis assumes $\volL$ accurately measures what
it claims.  The B-to-C gap is the scenario where this assumption
fails: $\volL$ diverges from the experiential volume it is supposed
to track.  A complete treatment would need to extend the phase
boundary to account for the divergence between $\volL$ and the
realized capability volume.  This extension is future work.

\paragraph{Cryptographic verification.}
\citet{lasser2026exo} addresses observation-channel integrity; this paper
analyzes the dynamical consequences of channel \emph{availability}.
The combination is: exogenous verification ensures channels are honest,
this paper ensures enough honest channels survive.  A joint treatment
where the phase boundary accounts for both integrity
failures~\citep{lasser2026exo} and availability failures (this paper)
would provide a stronger combined guarantee.

\paragraph{World-model verification and the public ledger.}
The self-correction rate $\rSeff$ depends on the information
environment.  In the minimal case (actor uses only its own observation
channels), $\rS$ is bounded by the channel sensitivity of
Definition~\ref{def:redundancy}.  If world-model predictions are
committed to the verification ledger of~\citet{lasser2026exo}, the
actor gains access to a faster correction channel: other agents whose
models disagree can flag discrepancies through the ledger's
verification protocol, providing corrections that are both direct
(rather than inferred from behavioral residuals) and exogenously
verified.  This increases $\rSeff$ without increasing $\rhomincross$
(the channel count is unchanged, but each channel's effective
correction rate improves), widening the self-correcting basin.

If world models are not merely committed but \emph{shareable} (other
agents can read and build on them), the ledger becomes a bootstrapping
mechanism: new agents enter the population with a calibrated world
model derived from the public record rather than from scratch
observation.  This addresses the initialization sensitivity of
Theorem~\ref{thm:redundancy} (the design criterion depends on
$\Lyap_{\max}$, the initial world-model error): a population with a
shared world-model ledger starts new agents with lower $\Lyap_{\max}$,
reducing the redundancy threshold for the self-correcting basin.

\subsection{Open questions}

\begin{enumerate}
\item \textbf{Nonlinear phase boundary.}  The state-dependent extension
  of the phase boundary, where $\rS$, $\rW$, and $\rsub$ are functions
  of $(P_t, \Wm_t)$ rather than constants, would produce tighter bounds
  at the cost of a switched-systems analysis with state-dependent
  switching signals.

\item \textbf{Characterization of absorbing states.}  We prove
  convergence to \emph{a} monopolar fixed point but do not characterize
  the full set of absorbing states.  Partial monopolarity (some
  dimensions monopolar, others not) may be stable under certain
  parameter regimes.

\item \textbf{Optimal redundancy allocation.}  Given a fixed budget of
  observation channels, how should they be allocated across dimensions
  and substrates to maximize the self-correcting basin?  This is a
  resource-allocation problem over the channel coverage matrix.

\item \textbf{Empirical calibration.}  The worked example uses
  pedagogical parameter values.  Calibrating $\rS$, $\rW$, $\rsub$,
  and $\dmax$ from operational data (e.g., from AI governance
  structures, human-oversight-of-AI dynamics, or historical
  institutional feedback loops) would convert the phase boundary from
  a qualitative characterization to a quantitative prediction.
\end{enumerate}
